\documentclass[a4paper]{article}
\usepackage{parskip}
\usepackage{lipsum}

\def\nterm {Spring}
\def\nyear {2026}
\def\ncourse {Probability -- Problems}

\newcommand{\convd}{\Rightarrow}
\newcommand{\convw}{\xrightarrow{w}}
\newcommand{\convp}{\xrightarrow{p}}

\input{../header}
\author{\nauthor\vspace{5pt}\\
Carnegie Mellon University}

\newcommand{\TODO}{\textcolor{red}{\textbf{*** TO-DO ***}}}

\begin{document}
\maketitle

\tableofcontents

\setcounter{section}{1}

\section{Laws of large numbers}

\setcounter{subsection}{1}

\subsection{Weak laws of large numbers}

\begin{ex}[Durrett 2.2.1]
  Let $X_1, X_2, \dots$ be uncorrelated random variables with
  finite expectations and
  \[
  i^{-1} \var X_i \to 0 \text{ as $i \to \infty$}.
  \]
  Let $S_n = X_1 + \dots + X_n$ and $\mu_n = E S_n / n$. Show that 
  $S_n / n - \mu_n \to 0$ in probability as $n \to \infty$.
\end{ex}


\begin{proof}
  WLOG assume $E X_n = 0$ for all $n$. Now we just 
  need to show that $S_n / n \to 0$ in probability. 
  Fix $\delta > 0$, we next show 
  \[
  P \left( \abs{\frac{S_n}{n}} > 2 \delta \right)
  \to 0
  \]
  
  Let $\epsilon > 0$ be arbitrary. 
  Since $i^{-1} \var X_i \to 0$ as $i \to \infty$,
  there exists $N$ such that $i \geq N$ implies 
  $i^{-1} \var X_i < \epsilon$. 
  Let $n \geq N$, we have 
  \begin{align*} 
    P \left( \abs{\frac{S_n}{n}} > 2 \delta \right)
    &= P \left( \frac{1}{n} \abs{\sumi^n X_i} > 2 \delta \right) \\
    &\leq P \left( \frac{1}{n} \sumi^N \abs{X_i} + 
    \frac{1}{n} \abs{\sum_{i=N+1}^{n} X_i} > 2 \delta \right) \\
    &\leq P \left( \frac{1}{n} \sumi^N \abs{X_i} > 
    \delta \right)
    + P \left( \frac{1}{n} \abs{\sum_{i=N + 1}^n X_i} > 
    \delta \right)
  \end{align*}
  For the first term, Chebyshev's inequality gives
  \begin{align*}
    P \left( \frac{1}{n} \sumi^N \abs{X_i} > \delta \right)
    &\leq \frac{1}{n \delta} \sumi^N E \abs{X_i} \to 0
  \end{align*}
  as $n \to \infty$.
  For the second term, Chebyshev's inequality gives
  \begin{align*}
    P \left( \frac{1}{n} \abs{\sum_{i=N+1}^n X_i} > \delta \right)
    \leq \frac{1}{\delta^2} \var \left( \frac{1}{n} \sum_{i=N + 1}^n X_i \right)
    = \frac{1}{n^2 \delta^2} \var \left( \sum_{i=N+1}^n X_i \right).
  \end{align*}
  Since $X_i$'s are uncorrelated, it follows that 
  \begin{align*}
    P \left( \frac{1}{n} \abs{\sum_{i=N+1}^n X_i} > \delta \right)
    &\leq \frac{1}{n^2 \delta^2} \sum_{i=N + 1}^n \var (X_i) 
    \leq \frac{1}{n \delta^2} \sum_{i=N + 1}^n \frac{\var X_i}{i}
    \leq \frac{(n - N) \epsilon}{n \delta^2} \leq \frac{\epsilon}{\delta^2}.
  \end{align*}
  Since $\epsilon > 0$ is arbitrary, this completes the proof.

\end{proof}

\begin{ex}[Durrett 2.2.2]
  The $L^2$ weak law generalizes immediately to certain dependent
  sequences. Suppose $E X_n$ and $E X_n X_m \leq r (n - m)$ for $m \leq n$ (no
  absolute value on the left-hand side!) with $r (k) \to 0$ as $k \to \infty$.
  Show that $(X_1 + \dots + X_n) / n \to 0$ in probability.
\end{ex}

\begin{proof}
  \TODO

\end{proof}

\begin{ex}[Durrett 2.2.5]
  Let $X_1, X_2, \dots$ be i.i.d.\ with 
  \[
  P(X_i > x) = \frac{e}{x \log x}.
  \]
  for $x \geq e$.
  Show that $E \abs{X_i} = \infty$, but there is a sequence of constants 
  $\mu_n \to \infty$ so that $S_n / n - \mu_n \to 0$ in probability.
\end{ex}

\begin{proof}
  Note that $P(X_i < e) = 0$. It follows that 
  \[
    E \abs{X_i} 
    = \int_e^\infty P(X_i > x) \d x 
    = [\log \log x]_e^\infty 
    = \infty.
  \]
  However, 
  \[
  x P(X_i > x) = \frac{e}{\log x} \to 0.
  \]
  Hence, by a previous theorem, if $\mu_n = E X_1 \ind_{(\abs{X_1} \leq n)}$ 
  then $S_n / n - \mu_n \to 0$. Note that $\mu_n < \infty$ but 
  $\mu_n \to E X_i = \infty$ as $n \to \infty$ 
  by monotone convergence theorem.

\end{proof}

\subsection{Borel-Cantelli lemmas}

\begin{ex}[Durrett 2.3.8]
Let $\seqinfn{A_n}$ be a sequence of independent events with $P(A_n) < 1$
for all $n$. Show that $P (\cupinfn A_n) = 1$ 
implies $\suminfn P(A_n) = \infty$ and hence $P(A_n \text{ i.o.}) = 1$.
\end{ex}

\begin{proof}
  Suppose for contradiction that $\suminfn P(A_n) = M < \infty$. 
  From independence, we have 
  \begin{align*}
    P \left( \capn^N A_n^c \right) 
    = \prod_{n=1}^N (1 - P(A_n)).
  \end{align*}
  As $\suminfn P(A_n)$ converges, $P(A_n) \to 0$ as $n \to 0$. 
  It follows that there exists $\alpha < 1$ 
  such that $P(A_n) \leq \alpha < 1$ for all $n$. 
  Once we select $\beta > 0$ such that $1 - \alpha \geq \exp (\beta \alpha)$,
  we obtain 
  \begin{align*}
    P \left( \capn^N A_n^c \right) 
    = \prod_{n=1}^N (1 - P(A_n))
    \geq \exp \left( - \beta \sumn^N P(A_n) \right)
  \end{align*}
  This implies 
  \begin{align*}
    P \left( \cupinfn A_n \right)
    = 1 - P \left( \capn^\infty A_n^c \right)
    \leq 1 - \exp \left( - \beta \sumn^\infty P(A_n) \right)
  \end{align*}
  Hence, $P (\cupinfn A_n) \leq 1 - \exp (- \beta M) < 1$, contradiction.

\end{proof}

\begin{ex}[Durrett 2.3.9]
  Show that if $P (A_n) \to 0$ and $\suminfn P (A_n^c \cap A_{n + 1}) < \infty$,
  then $P (A_n \text{ i.o.}) = 0$. Find an example of a sequence 
  $\seqinfn{A_n}$ to which the result above can
  be applied but the Borel-Cantelli lemma cannot.
\end{ex}

\begin{proof}
  For any $N, M$, we have 
  \[
  \bigcup_{n = N}^M A_n 
  = A_N \cup \bigcup_{n = N}^{M - 1} (A_n^c \cap A_{n + 1})
  \]
  This implies 
  \begin{align*}
    P \left( \bigcup_{n = N}^M A_n \right) 
    \leq P (A_N) + \sum_{n = N}^{M - 1} P (A_n^c \cap A_{n + 1}).
  \end{align*}
  Letting $M \to \infty$ shows
  $P (\bigcup_{n = N}^\infty A_n) \leq P(A_n) + \sum_{n = N}^\infty 
  P (A_n^c \cap A_{n + 1})$. 
  Now letting $N \to \infty$ shows $P (\limsup A_n) = 0$, as desired.
  
  For an example, define $(\Omega, \fcal, P)$ by letting $\Omega = [0, 1]$,
  $\fcal$ be the Borel sets, and $P$ be the Lebesgue measure. 
  For $n \geq 1$, define $A_n$ via 
  \[
  A_n = \left[ 1 - 2^{-n} - n^{-1} , 1 - 2^{-n} \right].
  \]
  Then $P(A_n) = n^{-1}$ but $P(A_n^c \cap A_{n + 1}) = 2^{-n - 1}$.


\end{proof}

\begin{ex}[Durrett 2.3.12]
  Let $X_1, X_2, \dots$ be a sequence of r.v.'s on $(\Omega, \fcal, P)$ 
  where $\Omega$ is a countable set and $\fcal$ consists of all subsets 
  of $\Omega$. Show that $X_n \to X$ in probability implies 
  $X_n \to X$ a.s.
\end{ex}

\begin{proof}
  \TODO

\end{proof}

\begin{ex}[Durrett 2.3.14]
  Let $X_1, X_2, \dots$ be independent. Show that $\sup X_n < \infty$ a.s. if
  and only if $\suminfn P(X_n > A) < \infty$ for some $A$.
\end{ex}

\begin{proof}
  \TODO

\end{proof}

\end{document}